\documentclass[12pt,a4paper]{article}

% ── Paquetes ────────────────────────────────────────────────────────────────
\usepackage[utf8]{inputenc}
\usepackage[T1]{fontenc}
\usepackage[spanish,mexico]{babel}
\usepackage{amsmath,amssymb,amsthm}
\usepackage{graphicx}
\usepackage{booktabs}
\usepackage{geometry}
\usepackage{hyperref}
\usepackage{listings}
\usepackage{xcolor}
\usepackage{float}
\usepackage{caption}
\usepackage{subcaption}
\usepackage{setspace}
\usepackage{parskip}
\usepackage{enumitem}
\usepackage{mathtools}
\usepackage{bm}
\usepackage{lmodern}
\usepackage{microtype}
\usepackage{fancyhdr}
\usepackage{titlesec}
\usepackage{tcolorbox}

\geometry{
  top=2.5cm, bottom=2.5cm,
  left=3cm,  right=3cm
}

% ── Estilo de página ─────────────────────────────────────────────────────────
\pagestyle{fancy}
\fancyhf{}
\rhead{\small Modelos SARIMA en Finanzas}
\lhead{\small Series de Tiempo}
\cfoot{\thepage}
\renewcommand{\headrulewidth}{0.4pt}

% ── Colores y listings ───────────────────────────────────────────────────────
\definecolor{codebg}{RGB}{245,245,245}
\definecolor{codeframe}{RGB}{180,180,180}
\definecolor{keycolor}{RGB}{0,80,160}
\definecolor{strcolor}{RGB}{160,0,0}
\definecolor{commentcolor}{RGB}{100,100,100}
\definecolor{numcolor}{RGB}{0,130,0}

\lstdefinestyle{pythonstyle}{
  language=Python,
  backgroundcolor=\color{codebg},
  basicstyle=\ttfamily\footnotesize,
  keywordstyle=\color{keycolor}\bfseries,
  stringstyle=\color{strcolor},
  commentstyle=\color{commentcolor}\itshape,
  numberstyle=\tiny\color{numcolor},
  frame=single,
  rulecolor=\color{codeframe},
  numbers=left,
  numbersep=6pt,
  breaklines=true,
  breakatwhitespace=false,
  tabsize=4,
  showstringspaces=false,
  captionpos=b,
  xleftmargin=12pt,
  xrightmargin=4pt,
  aboveskip=8pt,
  belowskip=8pt,
  morekeywords={import,from,as,def,return,with,True,False,None},
}

\lstset{style=pythonstyle}

% ── Entornos matemáticos ─────────────────────────────────────────────────────
\newtheorem{definicion}{Definición}[section]
\newtheorem{proposicion}{Proposición}[section]
\newtheorem{observacion}{Observación}[section]

% ── Operadores ───────────────────────────────────────────────────────────────
\DeclareMathOperator{\Cov}{Cov}
\DeclareMathOperator{\Var}{Var}
\DeclareMathOperator{\E}{\mathbb{E}}
\DeclareMathOperator{\Corr}{Corr}
\newcommand{\B}{\mathsf{B}}          % operador de rezago
\newcommand{\diff}{\nabla}           % operador diferencia

% ── Título ───────────────────────────────────────────────────────────────────
\title{
  \vspace{-1cm}
  {\Large \textbf{Modelos SARIMA en Series de Tiempo Financieras}} \\[0.6em]
  {\normalsize Teoría, Identificación, Estimación y Aplicaciones Prácticas con Python}
}
\author{
  \normalsize Notas de curso: Series de Tiempo Aplicadas a Finanzas \\
  \normalsize Universidad Panamericana --- Ingeniería Financiera
}
\date{\normalsize Abril 2026}

% ════════════════════════════════════════════════════════════════════════════
\begin{document}
% ════════════════════════════════════════════════════════════════════════════

\maketitle
\thispagestyle{fancy}
\tableofcontents
\newpage

% ────────────────────────────────────────────────────────────────────────────
\section{Introducción}
% ────────────────────────────────────────────────────────────────────────────

El análisis de series de tiempo es uno de los pilares fundamentales de las finanzas cuantitativas.
Desde la predicción de precios de activos y tipos de cambio hasta la modelación de volatilidad y
riesgo de crédito, los modelos paramétricos de series de tiempo permiten extraer estructura
estadística de datos secuenciales y generar pronósticos probabilísticos rigurosos.

Dentro de esta familia de modelos, los modelos \textbf{ARIMA} (\textit{Autoregressive Integrated
Moving Average}) propuestos formalmente por Box y Jenkins (1970) constituyen la base de la
econometría de series de tiempo univariadas. Su extensión natural, los modelos
\textbf{SARIMA} (\textit{Seasonal ARIMA}), incorpora patrones estacionales deterministas
presentes en muchas series financieras: rendimientos mensuales de fondos de inversión,
tasas de interés con ciclos anuales, volúmenes de negociación con patrones semanales, entre
otros.

\subsection{Relevancia en Finanzas}

Las aplicaciones de los modelos SARIMA en finanzas son amplias:

\begin{itemize}[leftmargin=*]
  \item \textbf{Tipos de cambio}: Las series de FX presentan con frecuencia autocorrelaciones
        de corto plazo y ocasionalmente patrones estacionales asociados a flujos comerciales.
  \item \textbf{Tasas de interés}: Las tasas de política monetaria y curvas de rendimiento
        exhiben alta persistencia y ciclos asociados a decisiones de bancos centrales.
  \item \textbf{Volumen bursátil}: Los volúmenes de negociación presentan estacionalidad
        intradiaria, semanal y mensual (días de vencimiento de opciones, rebalanceos de índices).
  \item \textbf{Commodities}: Los precios de energía y agrícolas presentan estacionalidad
        climática y de demanda.
  \item \textbf{Riesgo de crédito}: Los índices de morosidad pueden mostrar patrones
        estacionales relacionados con ciclos económicos.
\end{itemize}

\subsection{Limitaciones y contexto}

Es importante señalar que los modelos SARIMA pertenecen a la clase de modelos lineales para
la media condicional de la serie. No modelan directamente la volatilidad condicional (para
ello existen modelos GARCH), ni capturan no linealidades complejas (modelos de redes neuronales,
SETAR, etc.). Sin embargo, su interpretabilidad y parsimonia los hacen herramientas de referencia
obligatoria en cualquier análisis de series de tiempo.

% ────────────────────────────────────────────────────────────────────────────
\section{Forma General del Modelo SARIMA}
% ────────────────────────────────────────────────────────────────────────────

\subsection{Notación y operadores}

Denotamos la serie de tiempo como $\{y_t\}_{t=1}^{T}$, donde $t$ indexa el tiempo.
Introducimos dos operadores esenciales:

\begin{definicion}[Operador de rezago]
El \textbf{operador de rezago} $\B$ se define como:
\[
  \B^k y_t = y_{t-k}, \quad k \in \mathbb{Z}_{\geq 0}.
\]
\end{definicion}

\begin{definicion}[Operador de diferencia ordinaria]
El \textbf{operador de diferencia} de orden $d$ se define como:
\[
  \diff^d y_t = (1 - \B)^d y_t.
\]
En particular, $\diff y_t = y_t - y_{t-1}$ (diferencia de primer orden).
\end{definicion}

\begin{definicion}[Operador de diferencia estacional]
El \textbf{operador de diferencia estacional} de orden $D$ con periodo $s$ se define como:
\[
  \diff_s^D y_t = (1 - \B^s)^D y_t.
\]
Para $D=1$: $\diff_s y_t = y_t - y_{t-s}$.
\end{definicion}

\subsection{Definición formal del modelo SARIMA$(p,d,q)(P,D,Q)_s$}

\begin{definicion}[Modelo SARIMA]
Un proceso $\{y_t\}$ sigue un modelo SARIMA$(p,d,q)(P,D,Q)_s$ si satisface la ecuación:
\begin{equation}
  \label{eq:sarima}
  \boxed{
    \Phi_P(\B^s)\,\phi_p(\B)\,(1-\B)^d(1-\B^s)^D\,y_t
    = \Theta_Q(\B^s)\,\theta_q(\B)\,\varepsilon_t
  }
\end{equation}
donde $\varepsilon_t \sim \text{RB}(0,\sigma^2)$ (ruido blanco con varianza $\sigma^2$), y los
polinomios autorregresivos y de media móvil se definen como sigue.
\end{definicion}

\paragraph{Polinomios no estacionales:}
\begin{align}
  \phi_p(\B)   &= 1 - \phi_1 \B - \phi_2 \B^2 - \cdots - \phi_p \B^p, \label{eq:phi}\\
  \theta_q(\B) &= 1 + \theta_1 \B + \theta_2 \B^2 + \cdots + \theta_q \B^q. \label{eq:theta}
\end{align}

\paragraph{Polinomios estacionales:}
\begin{align}
  \Phi_P(\B^s)   &= 1 - \Phi_1 \B^s - \Phi_2 \B^{2s} - \cdots - \Phi_P \B^{Ps}, \label{eq:Phi}\\
  \Theta_Q(\B^s) &= 1 + \Theta_1 \B^s + \Theta_2 \B^{2s} + \cdots + \Theta_Q \B^{Qs}. \label{eq:Theta}
\end{align}

\subsection{Descripción de parámetros}

\subsubsection{Componentes no estacionales}

\begin{description}[leftmargin=3em,style=nextline]
  \item[$p$ --- Orden autorregresivo]
    Indica cuántos rezagos de la serie (ya diferenciada) entran como predictores lineales.
    Un AR($p$) captura la inercia o memoria de corto plazo del proceso.
    \textit{Interpretación financiera}: $p=1$ en rendimientos implica que el rendimiento
    de hoy está parcialmente determinado por el de ayer (posible autocorrelación de primer orden,
    frecuente en retornos diarios de activos ilíquidos o con efectos momentum débiles).

  \item[$d$ --- Orden de integración ordinaria]
    Número de diferencias ordinarias necesarias para que la serie sea estacionaria.
    En series de precios de activos financieros se tiene típicamente $d=1$, pues los precios
    son no estacionarios (caminata aleatoria) pero los rendimientos logarítmicos son estacionarios.
    La condición de estacionariedad requiere que todas las raíces de $\phi_p(z) = 0$ estén
    fuera del círculo unitario.

  \item[$q$ --- Orden de media móvil]
    Número de shocks pasados (errores de pronóstico) que afectan el nivel actual de la serie.
    Un MA($q$) implica que el proceso tiene \textit{memoria finita} de exactamente $q$ periodos.
    \textit{Interpretación financiera}: $q=1$ puede reflejar el ajuste de precios ante nueva
    información con un rezago (e.g., mercados con liquidez reducida).
\end{description}

\subsubsection{Componentes estacionales}

\begin{description}[leftmargin=3em,style=nextline]
  \item[$P$ --- Orden autorregresivo estacional]
    Número de rezagos estacionales de la serie diferenciada que actúan como predictores.
    Captura la dependencia entre observaciones separadas exactamente $s$ periodos.

  \item[$D$ --- Orden de integración estacional]
    Número de diferencias estacionales necesarias para eliminar tendencias estacionales.
    Una diferencia estacional $D=1$ sustrae el valor del mismo periodo del año (o semana, mes)
    anterior: $y_t - y_{t-s}$.

  \item[$Q$ --- Orden de media móvil estacional]
    Número de shocks estacionales pasados incluidos en el modelo.

  \item[$s$ --- Periodo estacional]
    Frecuencia del ciclo estacional. Valores comunes en finanzas:
    \begin{itemize}
      \item $s = 4$: datos trimestrales con ciclo anual.
      \item $s = 12$: datos mensuales con ciclo anual.
      \item $s = 5$ o $s=22$: datos diarios con ciclo semanal o mensual.
      \item $s = 252$: datos diarios con ciclo anual (días hábiles).
    \end{itemize}
\end{description}

\subsection{Condiciones de estacionariedad e invertibilidad}

Para que el modelo SARIMA sea estacionario (en la parte no integrada) e invertible, se
requiere que todas las raíces de los polinomios $\phi_p(z) = 0$ y $\Phi_P(z^s) = 0$ estén
fuera del círculo unitario en el plano complejo, y que análogamente todas las raíces de
$\theta_q(z) = 0$ y $\Theta_Q(z^s) = 0$ estén fuera del círculo unitario.

% ────────────────────────────────────────────────────────────────────────────
\section{Herramientas de Identificación}
% ────────────────────────────────────────────────────────────────────────────

\subsection{Función de Autocorrelación (ACF)}

\begin{definicion}[ACF]
Para un proceso estacionario $\{y_t\}$ con media $\mu$ y varianza $\sigma_y^2 < \infty$,
la \textbf{función de autocorrelación} (ACF) en el rezago $k$ se define como:
\[
  \rho(k) = \frac{\Cov(y_t, y_{t-k})}{\Var(y_t)}
           = \frac{\gamma(k)}{\gamma(0)}, \quad k = 0, \pm 1, \pm 2, \ldots
\]
donde $\gamma(k) = \Cov(y_t, y_{t-k})$ es la función de autocovarianza.
\end{definicion}

\paragraph{Intuición.}
La ACF mide la correlación lineal entre $y_t$ y su propio pasado $y_{t-k}$, sin controlar
por los rezagos intermedios. En la muestra, se estima como:
\[
  \hat{\rho}(k) = \frac{\displaystyle\sum_{t=k+1}^{T}(y_t - \bar{y})(y_{t-k} - \bar{y})}
                       {\displaystyle\sum_{t=1}^{T}(y_t - \bar{y})^2}.
\]

\paragraph{Uso en selección de modelos.}
\begin{itemize}
  \item Un proceso \textbf{MA($q$)} presenta ACF que se \textit{trunca} exactamente en el rezago $q$:
        $\rho(k) = 0$ para $|k| > q$.
  \item Un proceso \textbf{AR($p$)} presenta ACF que \textit{decae} exponencialmente (o de forma
        sinusoidal amortiguada), sin truncamiento abrupto.
  \item La presencia de picos significativos en los rezagos $k = s, 2s, 3s, \ldots$ indica
        componentes estacionales.
  \item Las bandas de confianza aproximadas al $95\%$ son $\pm 1.96 / \sqrt{T}$.
\end{itemize}

\subsection{Función de Autocorrelación Parcial (PACF)}

\begin{definicion}[PACF]
La \textbf{función de autocorrelación parcial} (PACF) en el rezago $k$, denotada $\alpha(k)$,
mide la correlación entre $y_t$ e $y_{t-k}$ \emph{eliminando} el efecto lineal de las variables
intermedias $y_{t-1}, y_{t-2}, \ldots, y_{t-k+1}$:
\[
  \alpha(k) = \Corr(y_t - \hat{y}_t^{(k)},\; y_{t-k} - \hat{y}_{t-k}^{(k)}),
\]
donde $\hat{y}_t^{(k)}$ es la proyección lineal de $y_t$ sobre $\{y_{t-1}, \ldots, y_{t-k+1}\}$.
\end{definicion}

\paragraph{Intuición.}
La PACF responde a la pregunta: ¿después de descontar la influencia de los rezagos
$1, 2, \ldots, k-1$, aún existe correlación significativa entre $y_t$ e $y_{t-k}$?

\paragraph{Uso en selección de modelos.}
\begin{itemize}
  \item Un proceso \textbf{AR($p$)} presenta PACF que se \textit{trunca} exactamente en el
        rezago $p$: $\alpha(k) = 0$ para $k > p$.
  \item Un proceso \textbf{MA($q$)} presenta PACF que \textit{decae} gradualmente.
  \item Combinando ACF y PACF se puede identificar tentativamente el orden $(p,q)$ del modelo.
\end{itemize}

\begin{table}[H]
\centering
\caption{Patrones de ACF y PACF según el proceso generador de datos}
\label{tab:acf_pacf}
\begin{tabular}{lccc}
\toprule
\textbf{Proceso} & \textbf{ACF} & \textbf{PACF} \\
\midrule
AR($p$) & Decae exponencialmente & Se trunca en rezago $p$ \\
MA($q$) & Se trunca en rezago $q$ & Decae gradualmente \\
ARMA($p,q$) & Decae gradualmente & Decae gradualmente \\
RB (Ruido blanco) & $\approx 0$ para $k \geq 1$ & $\approx 0$ para $k \geq 1$ \\
\bottomrule
\end{tabular}
\end{table}

\subsection{Criterio de Información de Akaike (AIC)}

\begin{definicion}[AIC]
Sea $\hat{L}$ el valor máximo de la función de verosimilitud del modelo con $k$ parámetros
libres. El \textbf{criterio de Akaike} se define como:
\[
  \mathrm{AIC} = -2\ln(\hat{L}) + 2k.
\]
\end{definicion}

\paragraph{Intuición.}
El AIC penaliza la complejidad del modelo (número de parámetros $k$) frente al ajuste
(log-verosimilitud). El término $-2\ln(\hat{L})$ disminuye al añadir parámetros (mejor ajuste),
pero el término $2k$ crece linealmente. El AIC busca el \textit{balance óptimo}.

\paragraph{Uso práctico.}
Se prefiere el modelo con \textbf{menor AIC}. Al comparar modelos SARIMA con distintos
órdenes $(p,q,P,Q)$, el AIC tiende a seleccionar modelos más parsimoniosos pero puede
\textit{sobreajustar} ligeramente en muestras pequeñas.

\subsection{Criterio de Información Bayesiano (BIC)}

\begin{definicion}[BIC]
\[
  \mathrm{BIC} = -2\ln(\hat{L}) + k\ln(T),
\]
donde $T$ es el tamaño de la muestra.
\end{definicion}

\paragraph{Intuición y comparación con AIC.}
El BIC penaliza más severamente la complejidad que el AIC (pues $\ln(T) > 2$ para $T > 7$).
Como consecuencia, el BIC es \textit{consistente}: cuando $T \to \infty$, selecciona el verdadero
modelo generador de datos (si este pertenece a la familia considerada), mientras que el AIC
tiende a elegir modelos ligeramente sobreajustados.

\paragraph{Regla práctica.}
\begin{itemize}
  \item Si AIC y BIC coinciden en el modelo seleccionado, existe evidencia sólida.
  \item Si difieren, se prefiere BIC en muestras grandes y AIC en muestras pequeñas o cuando
        el objetivo es minimizar el error de pronóstico (no identificar el verdadero modelo).
\end{itemize}

\subsection{Diferenciación: $d$ y $D$}

La diferenciación tiene como objetivo transformar una serie no estacionaria en una serie
estacionaria. Las pruebas formales para determinar $d$ y $D$ incluyen:

\begin{itemize}
  \item \textbf{Prueba de Dickey-Fuller Aumentada (ADF)}: contrasta $H_0$: raíz unitaria vs.
        $H_1$: estacionariedad.
  \item \textbf{Prueba KPSS}: contrasta $H_0$: estacionariedad vs. $H_1$: raíz unitaria.
  \item \textbf{Criterio visual}: gráfica de la serie y correlograma (ACF muy persistente
        sugiere no estacionariedad).
\end{itemize}

Para la diferenciación estacional, se aplica $\diff_s^D y_t = (1 - \B^s)^D y_t$ y se verifica
que los picos estacionales desaparecen del ACF.

\subsection{Estacionariedad}

\begin{definicion}[Estacionariedad débil]
Un proceso $\{y_t\}$ es \textbf{débilmente estacionario} (o estacionario de covarianza) si:
\begin{enumerate}
  \item $\E[y_t] = \mu < \infty$ para todo $t$ (media constante),
  \item $\Var(y_t) = \sigma^2 < \infty$ para todo $t$ (varianza finita y constante),
  \item $\Cov(y_t, y_{t-k}) = \gamma(k)$ depende solo de $k$, no de $t$ (autocovarianzas
        independientes del tiempo).
\end{enumerate}
\end{definicion}

\paragraph{Relevancia en finanzas.}
Los precios de activos son típicamente no estacionarios (procesos integrados de orden 1,
$I(1)$), mientras que los rendimientos logarítmicos $r_t = \ln(P_t/P_{t-1})$ suelen ser
estacionarios $I(0)$. Trabajar con la serie no estacionaria directamente llevaría a
\textit{regresiones espurias} con $R^2$ elevados y estadísticos $t$ no confiables.

\subsection{Ruido Blanco}

\begin{definicion}[Ruido blanco]
Un proceso $\{\varepsilon_t\}$ es \textbf{ruido blanco} si:
\[
  \E[\varepsilon_t] = 0, \quad \E[\varepsilon_t^2] = \sigma^2, \quad \E[\varepsilon_t \varepsilon_s] = 0 \text{ para } t \neq s.
\]
Si adicionalmente $\varepsilon_t \sim \mathcal{N}(0, \sigma^2)$, se denomina \textbf{ruido blanco gaussiano}.
\end{definicion}

\paragraph{Uso diagnóstico.}
Los residuos de un modelo SARIMA bien especificado deben comportarse como ruido blanco.
Se verifica mediante:
\begin{itemize}
  \item \textbf{Prueba de Ljung-Box}: contrasta $H_0$: no autocorrelación en los residuos
        hasta el rezago $m$.
        \[
          Q_{LB} = T(T+2)\sum_{k=1}^{m}\frac{\hat{\rho}_\varepsilon^2(k)}{T-k} \overset{H_0}{\sim} \chi^2(m - p - q - P - Q).
        \]
  \item \textbf{Prueba de Jarque-Bera}: verifica normalidad de los residuos.
\end{itemize}

% ────────────────────────────────────────────────────────────────────────────
\section{Metodología Box-Jenkins}
% ────────────────────────────────────────────────────────────────────────────

La metodología Box-Jenkins es un procedimiento iterativo de cuatro etapas para la
construcción de modelos SARIMA:

\subsection{Etapa 1: Identificación}

\begin{enumerate}
  \item \textbf{Graficar la serie}: detectar visualmente tendencias, quiebres estructurales,
        heterocedasticidad y patrones estacionales.
  \item \textbf{Transformaciones}: si la varianza es no constante (heterocedasticidad),
        aplicar la transformación logarítmica $\tilde{y}_t = \ln(y_t)$ o la transformación
        Box-Cox:
        \[
          \tilde{y}_t = \begin{cases} (y_t^\lambda - 1)/\lambda & \lambda \neq 0 \\ \ln(y_t) & \lambda = 0 \end{cases}.
        \]
  \item \textbf{Diferenciación}: determinar $d$ y $D$ mediante pruebas de raíz unitaria
        (ADF, KPSS) y análisis visual del ACF.
  \item \textbf{Análisis ACF y PACF}: sobre la serie diferenciada estacionaria, examinar
        los correlogramas para proponer órdenes tentativos $(p, q)$ y $(P, Q)$.
  \item \textbf{Definir candidatos}: listar un conjunto de modelos candidatos
        SARIMA$(p,d,q)(P,D,Q)_s$ con orden bajo (parsimonia).
\end{enumerate}

\subsection{Etapa 2: Estimación}

Los parámetros del modelo SARIMA se estiman típicamente por \textbf{máxima verosimilitud
condicional} (MLE). Dado el modelo~\eqref{eq:sarima}, la función de log-verosimilitud
condicional bajo normalidad de $\varepsilon_t$ es:
\[
  \ell(\bm{\phi}, \bm{\theta}, \bm{\Phi}, \bm{\Theta}, \sigma^2 \mid \mathbf{y})
  = -\frac{T^*}{2}\ln(2\pi\sigma^2) - \frac{1}{2\sigma^2}\sum_{t}\varepsilon_t^2,
\]
donde $T^* = T - d - Ds$ es el número efectivo de observaciones tras la diferenciación, y
$\varepsilon_t$ son los residuos generados recursivamente por el modelo.

Los estimadores de MLE son asintóticamente normales bajo condiciones de regularidad, lo que
permite construir intervalos de confianza y pruebas de hipótesis estándar.

\subsection{Etapa 3: Validación (selección de modelo)}

Una vez estimados los modelos candidatos:
\begin{enumerate}
  \item Comparar AIC y BIC entre modelos: seleccionar el de menor valor.
  \item Verificar que todos los coeficientes estimados sean estadísticamente significativos
        (valor-$p < 0.05$ o intervalos de confianza que excluyan cero).
  \item Comprobar la invertibilidad y estacionariedad: las raíces de los polinomios
        estimados deben estar fuera del círculo unitario.
  \item Principio de parsimonia: ante modelos con AIC/BIC similares, preferir el más simple.
\end{enumerate}

\subsection{Etapa 4: Diagnóstico de residuos}

El modelo seleccionado es adecuado si los residuos $\hat{\varepsilon}_t$ se comportan
como ruido blanco:
\begin{enumerate}
  \item \textbf{Gráfica de residuos}: no debe haber patrones, tendencias ni heterocedasticidad visible.
  \item \textbf{ACF/PACF de residuos}: todos los picos deben estar dentro de las bandas
        de confianza al 95\%.
  \item \textbf{Prueba de Ljung-Box}: $p$-valor$> 0.05$ indica ausencia de autocorrelación.
  \item \textbf{Normalidad}: histograma y QQ-plot de residuos; prueba de Jarque-Bera.
  \item Si el diagnóstico falla, se regresa a la Etapa 1 y se reidentifica el modelo.
\end{enumerate}

% ────────────────────────────────────────────────────────────────────────────
\section{Ejemplos Prácticos con Python}
% ────────────────────────────────────────────────────────────────────────────

Los tres ejemplos siguientes utilizan datos financieros reales descargados mediante
\texttt{yfinance}. El paquete \texttt{statsmodels} provee la implementación de SARIMA
a través de \texttt{SARIMAX}. Se asume la instalación previa de:
\begin{center}
  \texttt{pip install yfinance statsmodels pandas matplotlib scipy}
\end{center}

% ──────────────────────────────────────────
\subsection{Ejemplo 1: Tipo de cambio USD/MXN --- Diferenciación y ARIMA}
% ──────────────────────────────────────────

\subsubsection{Descripción de los datos}

El tipo de cambio peso mexicano frente al dólar estadounidense (USD/MXN) es una serie
diaria con alta persistencia. Los precios del FX son $I(1)$: la serie de niveles tiene
raíz unitaria, pero los rendimientos logarítmicos diarios son estacionarios. Utilizaremos
datos de los últimos dos años para ajustar un modelo ARIMA a los rendimientos.

\subsubsection{Código Python}

\begin{lstlisting}[caption={Ejemplo 1: ARIMA para rendimientos USD/MXN},label=lst:usdmxn]
import yfinance as yf
import pandas as pd
import numpy as np
import matplotlib.pyplot as plt
from statsmodels.tsa.stattools import adfuller, acf, pacf
from statsmodels.graphics.tsaplots import plot_acf, plot_pacf
from statsmodels.tsa.statespace.sarimax import SARIMAX
from scipy import stats
import warnings
warnings.filterwarnings('ignore')

# ── 1. Descarga de datos ────────────────────────────────────
ticker = yf.Ticker("MXN=X")
df = ticker.history(period="2y", interval="1d")
precio = df['Close'].dropna()
rendimiento = np.log(precio).diff().dropna()

# ── 2. Visualizacion ────────────────────────────────────────
fig, axes = plt.subplots(2, 1, figsize=(12, 6), sharex=False)
axes[0].plot(precio, color='steelblue', linewidth=0.8)
axes[0].set_title('Precio USD/MXN (cierre diario)')
axes[0].set_ylabel('Pesos por dolar')
axes[1].plot(rendimiento, color='darkred', linewidth=0.6)
axes[1].axhline(0, color='black', linewidth=0.5)
axes[1].set_title('Rendimiento logaritmico diario')
axes[1].set_ylabel('log-retorno')
plt.tight_layout()
plt.savefig('usdmxn_serie.png', dpi=150, bbox_inches='tight')
plt.show()

# ── 3. Prueba ADF ───────────────────────────────────────────
def adf_test(serie, nombre):
    resultado = adfuller(serie, autolag='AIC')
    print(f"\n=== ADF: {nombre} ===")
    print(f"  Estadistico ADF : {resultado[0]:.4f}")
    print(f"  p-valor         : {resultado[1]:.4f}")
    print(f"  Valores criticos: {resultado[4]}")
    if resultado[1] < 0.05:
        print("  -> Rechazamos H0: la serie ES estacionaria.")
    else:
        print("  -> No rechazamos H0: la serie NO es estacionaria.")

adf_test(precio, "Precio USD/MXN")
adf_test(rendimiento, "Rendimiento log USD/MXN")

# ── 4. ACF y PACF ───────────────────────────────────────────
fig, axes = plt.subplots(1, 2, figsize=(12, 4))
plot_acf(rendimiento, lags=30, ax=axes[0], title='ACF - Rendimiento USD/MXN')
plot_pacf(rendimiento, lags=30, ax=axes[1], title='PACF - Rendimiento USD/MXN')
plt.tight_layout()
plt.savefig('usdmxn_acf_pacf.png', dpi=150, bbox_inches='tight')
plt.show()

# ── 5. Seleccion de modelo (grid search AIC) ────────────────
resultados = []
for p in range(0, 4):
    for q in range(0, 4):
        try:
            modelo = SARIMAX(rendimiento, order=(p, 0, q),
                             trend='n', enforce_stationarity=True,
                             enforce_invertibility=True)
            res = modelo.fit(disp=False)
            resultados.append({'p': p, 'q': q,
                                'AIC': res.aic, 'BIC': res.bic})
        except Exception:
            pass

df_res = pd.DataFrame(resultados).sort_values('AIC')
print("\n=== Top 5 modelos por AIC ===")
print(df_res.head())

# ── 6. Modelo optimo ────────────────────────────────────────
mejor = df_res.iloc[0]
p_opt, q_opt = int(mejor['p']), int(mejor['q'])
print(f"\nModelo seleccionado: ARIMA({p_opt}, 0, {q_opt})")

modelo_final = SARIMAX(rendimiento, order=(p_opt, 0, q_opt),
                       trend='n').fit(disp=False)
print(modelo_final.summary())

# ── 7. Diagnostico de residuos ──────────────────────────────
residuos = modelo_final.resid
fig, axes = plt.subplots(2, 2, figsize=(12, 8))
axes[0, 0].plot(residuos, color='navy', linewidth=0.6)
axes[0, 0].set_title('Residuos')
plot_acf(residuos, lags=20, ax=axes[0, 1], title='ACF Residuos')
axes[1, 0].hist(residuos, bins=40, color='steelblue', edgecolor='white')
axes[1, 0].set_title('Histograma de Residuos')
stats.probplot(residuos, dist='norm', plot=axes[1, 1])
axes[1, 1].set_title('QQ-Plot')
plt.tight_layout()
plt.savefig('usdmxn_diagnostico.png', dpi=150, bbox_inches='tight')
plt.show()

# Prueba Ljung-Box
from statsmodels.stats.diagnostic import acorr_ljungbox
lb = acorr_ljungbox(residuos, lags=[10, 20], return_df=True)
print("\n=== Prueba Ljung-Box ===")
print(lb)

# ── 8. Prediccion ───────────────────────────────────────────
n_pred = 10
pred = modelo_final.get_forecast(steps=n_pred)
pred_ci = pred.conf_int(alpha=0.05)
pred_mean = pred.predicted_mean

print("\n=== Prediccion 10 dias adelante (rendimiento) ===")
print(pd.DataFrame({'Pred': pred_mean,
                    'IC_inf': pred_ci.iloc[:, 0],
                    'IC_sup': pred_ci.iloc[:, 1]}))
\end{lstlisting}

\subsubsection{Interpretación de resultados}

La prueba ADF típicamente rechaza la raíz unitaria en los rendimientos logarítmicos del
USD/MXN, confirmando la estacionariedad. El análisis ACF/PACF de los rendimientos suele
mostrar autocorrelaciones pequeñas pero significativas en algunos rezagos, lo que motiva
la inclusión de términos AR y/o MA de orden bajo. El grid search selecciona el modelo
con menor AIC, generalmente ARIMA(1,0,1) o ARIMA(0,0,1) para rendimientos de FX.
El diagnóstico de residuos verifica que los errores del modelo no presentan estructura
remanente: la prueba de Ljung-Box no debe rechazar $H_0$ de no autocorrelación.

% ──────────────────────────────────────────
\subsection{Ejemplo 2: Retornos mensuales del S\&P 500 --- SARIMA con componente estacional}
% ──────────────────────────────────────────

\subsubsection{Descripción de los datos}

El índice S\&P 500 (ticker \texttt{GSPC}) presentado en frecuencia mensual puede exhibir
patrones estacionales relacionados con ciclos de ganancias corporativas, rebalanceos
institucionales y efectos de calendario (``enero'', ``venta en mayo'', etc.). Utilizaremos
10 años de datos mensuales y ajustaremos un SARIMA$(p,d,q)(P,D,Q)_{12}$.

\subsubsection{Código Python}

\begin{lstlisting}[caption={Ejemplo 2: SARIMA para S\&P 500 mensual},label=lst:sp500]
import yfinance as yf
import pandas as pd
import numpy as np
import matplotlib.pyplot as plt
from statsmodels.tsa.stattools import adfuller
from statsmodels.graphics.tsaplots import plot_acf, plot_pacf
from statsmodels.tsa.statespace.sarimax import SARIMAX
import warnings
warnings.filterwarnings('ignore')

# ── 1. Descarga y resampleo mensual ────────────────────────
sp500 = yf.download("^GSPC", period="10y", interval="1mo", progress=False)
precio_m = sp500['Close'].dropna()
retorno_m = np.log(precio_m).diff().dropna() * 100   # en porcentaje

print(f"Observaciones: {len(retorno_m)}")
print(retorno_m.describe())

# ── 2. Visualizacion ────────────────────────────────────────
fig, axes = plt.subplots(2, 1, figsize=(12, 6))
axes[0].plot(precio_m, color='steelblue')
axes[0].set_title('S&P 500 (precio mensual de cierre)')
axes[1].bar(retorno_m.index, retorno_m, color=np.where(retorno_m >= 0,
            'green', 'red'), width=20)
axes[1].set_title('Retorno logaritmico mensual (%)')
plt.tight_layout()
plt.savefig('sp500_serie.png', dpi=150, bbox_inches='tight')
plt.show()

# ── 3. Prueba ADF ───────────────────────────────────────────
r_adf = adfuller(retorno_m, autolag='AIC')
print(f"\nADF retornos mensuales: stat={r_adf[0]:.4f}, p={r_adf[1]:.4f}")

# ── 4. ACF/PACF (hasta rezago 36) ──────────────────────────
fig, axes = plt.subplots(1, 2, figsize=(14, 4))
plot_acf(retorno_m, lags=36, ax=axes[0],
         title='ACF - Retorno mensual S&P 500')
plot_pacf(retorno_m, lags=36, ax=axes[1],
          title='PACF - Retorno mensual S&P 500')
plt.tight_layout()
plt.savefig('sp500_acf_pacf.png', dpi=150, bbox_inches='tight')
plt.show()

# ── 5. Grid search SARIMA ────────────────────────────────────
S = 12    # estacionalidad mensual
resultados = []
for p in range(0, 3):
    for q in range(0, 3):
        for P in range(0, 2):
            for Q in range(0, 2):
                try:
                    m = SARIMAX(retorno_m,
                                order=(p, 0, q),
                                seasonal_order=(P, 0, Q, S),
                                trend='c',
                                enforce_stationarity=True,
                                enforce_invertibility=True)
                    res = m.fit(disp=False, maxiter=200)
                    resultados.append({
                        'p': p, 'q': q, 'P': P, 'Q': Q,
                        'AIC': res.aic, 'BIC': res.bic
                    })
                except Exception:
                    pass

df_grid = pd.DataFrame(resultados).sort_values('AIC').reset_index(drop=True)
print("\n=== Top 5 modelos SARIMA por AIC ===")
print(df_grid.head())

# ── 6. Ajuste del modelo optimo ─────────────────────────────
fila = df_grid.iloc[0]
p_o, q_o, P_o, Q_o = int(fila.p), int(fila.q), int(fila.P), int(fila.Q)
print(f"\nModelo: SARIMA({p_o},0,{q_o})({P_o},0,{Q_o})_12")

modelo = SARIMAX(retorno_m,
                 order=(p_o, 0, q_o),
                 seasonal_order=(P_o, 0, Q_o, S),
                 trend='c').fit(disp=False, maxiter=300)
print(modelo.summary())

# ── 7. Diagnostico ──────────────────────────────────────────
modelo.plot_diagnostics(figsize=(12, 8))
plt.suptitle('Diagnostico de residuos - S&P 500 SARIMA', y=1.02)
plt.tight_layout()
plt.savefig('sp500_diagnostico.png', dpi=150, bbox_inches='tight')
plt.show()

from statsmodels.stats.diagnostic import acorr_ljungbox
lb = acorr_ljungbox(modelo.resid, lags=[12, 24], return_df=True)
print("\n=== Ljung-Box ===")
print(lb)

# ── 8. Prediccion 12 meses ──────────────────────────────────
forecast = modelo.get_forecast(steps=12)
fc_mean = forecast.predicted_mean
fc_ci   = forecast.conf_int(alpha=0.05)

fig, ax = plt.subplots(figsize=(12, 5))
ax.plot(retorno_m[-36:], label='Historico', color='steelblue')
ax.plot(fc_mean, label='Prediccion', color='darkorange', linestyle='--')
ax.fill_between(fc_ci.index,
                fc_ci.iloc[:, 0], fc_ci.iloc[:, 1],
                alpha=0.25, color='orange', label='IC 95%')
ax.axhline(0, color='black', linewidth=0.5)
ax.set_title('Prediccion 12 meses - Retorno S&P 500')
ax.legend()
plt.tight_layout()
plt.savefig('sp500_forecast.png', dpi=150, bbox_inches='tight')
plt.show()
\end{lstlisting}

\subsubsection{Interpretación de resultados}

Los retornos logarítmicos mensuales del S\&P 500 son estacionarios (la ADF rechaza la raíz
unitaria), con media positiva pequeña (prima de riesgo) y desviación estándar elevada.
La ACF mensual puede mostrar picos en rezagos múltiplos de 12 (efecto enero, ciclos anuales),
lo que justifica la componente estacional con $s=12$. El BIC típicamente selecciona modelos
parsimoniosos como SARIMA$(1,0,0)(0,0,1)_{12}$ o similares. La predicción puntual a 12 meses
tiene intervalos de confianza amplios, lo cual es coherente con la alta incertidumbre inherente
a los retornos bursátiles.

% ──────────────────────────────────────────
\subsection{Ejemplo 3: Petróleo WTI --- SARIMA con diferenciación estacional}
% ──────────────────────────────────────────

\subsubsection{Descripción de los datos}

Los precios del petróleo crudo WTI (West Texas Intermediate) exhiben fuerte estacionalidad
relacionada con la demanda de combustible (verano en EE.UU., invierno en Europa) y los ciclos
de producción de la OPEP+. Trabajaremos con datos semanales de los últimos 5 años,
aplicando $d=1$ (precio no estacionario) y $D=1$ con $s=52$ (ciclo anual en semanas).

\subsubsection{Código Python}

\begin{lstlisting}[caption={Ejemplo 3: SARIMA para petróleo WTI semanal},label=lst:wti]
import yfinance as yf
import pandas as pd
import numpy as np
import matplotlib.pyplot as plt
from statsmodels.tsa.stattools import adfuller, kpss
from statsmodels.graphics.tsaplots import plot_acf, plot_pacf
from statsmodels.tsa.statespace.sarimax import SARIMAX
from statsmodels.stats.diagnostic import acorr_ljungbox
from scipy import stats
import warnings
warnings.filterwarnings('ignore')

# ── 1. Descarga datos WTI semanales ────────────────────────
wti = yf.download("CL=F", period="5y", interval="1wk", progress=False)
precio_wti = wti['Close'].dropna()
log_precio  = np.log(precio_wti)

print(f"Observaciones: {len(precio_wti)}")
print(precio_wti.describe())

# ── 2. Visualizacion ────────────────────────────────────────
fig, axes = plt.subplots(3, 1, figsize=(12, 9))
axes[0].plot(precio_wti, color='saddlebrown', linewidth=0.8)
axes[0].set_title('Precio WTI (cierre semanal, USD/barril)')
axes[1].plot(log_precio, color='darkorange', linewidth=0.8)
axes[1].set_title('Log-precio WTI')
# primera diferencia del log-precio
diff1 = log_precio.diff().dropna()
axes[2].plot(diff1, color='navy', linewidth=0.6)
axes[2].axhline(0, color='black', lw=0.5)
axes[2].set_title('Primera diferencia del log-precio WTI')
plt.tight_layout()
plt.savefig('wti_serie.png', dpi=150, bbox_inches='tight')
plt.show()

# ── 3. Pruebas de estacionariedad ───────────────────────────
def adf_kpss(serie, nombre):
    adf_r = adfuller(serie, autolag='AIC')
    kpss_r = kpss(serie, regression='c', nlags='auto')
    print(f"\n--- {nombre} ---")
    print(f"  ADF  stat={adf_r[0]:.4f}, p={adf_r[1]:.4f}  "
          f"({'estacionaria' if adf_r[1] < 0.05 else 'NO estacionaria'})")
    print(f"  KPSS stat={kpss_r[0]:.4f}, p={kpss_r[1]:.4f}  "
          f"({'NO estacionaria' if kpss_r[1] < 0.05 else 'estacionaria'})")

adf_kpss(log_precio, "Log-precio WTI (nivel)")
adf_kpss(diff1, "Diferencia log-precio WTI")

# ── 4. ACF y PACF de la serie diferenciada ──────────────────
fig, axes = plt.subplots(1, 2, figsize=(14, 4))
plot_acf(diff1, lags=60, ax=axes[0],
         title='ACF - Diff Log-precio WTI')
plot_pacf(diff1, lags=60, ax=axes[1],
          title='PACF - Diff Log-precio WTI')
plt.tight_layout()
plt.savefig('wti_acf_pacf.png', dpi=150, bbox_inches='tight')
plt.show()

# ── 5. Diferenciacion estacional (D=1, s=52) ────────────────
diff_estacional = log_precio.diff(52).diff(1).dropna()
print(f"\nObservaciones tras doble diferenciacion: {len(diff_estacional)}")

adf_kpss(diff_estacional, "Diff estacional + ordinaria")

fig, axes = plt.subplots(1, 2, figsize=(14, 4))
plot_acf(diff_estacional, lags=60, ax=axes[0],
         title='ACF - Diff estacional WTI')
plot_pacf(diff_estacional, lags=60, ax=axes[1],
          title='PACF - Diff estacional WTI')
plt.tight_layout()
plt.savefig('wti_acf_pacf_diff.png', dpi=150, bbox_inches='tight')
plt.show()

# ── 6. Modelos candidatos SARIMA ────────────────────────────
# Usamos s=52 (ciclo anual en semanas)
# Para reducir tiempo de computo, limitamos el espacio de busqueda
S_wti = 52
candidatos = [
    (1, 1, 0, 0, 1, 0),
    (0, 1, 1, 0, 1, 1),
    (1, 1, 1, 0, 1, 0),
    (1, 1, 1, 0, 1, 1),
    (2, 1, 1, 0, 1, 0),
    (1, 1, 2, 0, 1, 1),
    (0, 1, 1, 1, 1, 0),
    (1, 1, 0, 1, 1, 1),
]

resultados_wti = []
for (p, d, q, P, D, Q) in candidatos:
    try:
        m = SARIMAX(log_precio,
                    order=(p, d, q),
                    seasonal_order=(P, D, Q, S_wti),
                    trend='n',
                    enforce_stationarity=True,
                    enforce_invertibility=True)
        res = m.fit(disp=False, maxiter=300)
        resultados_wti.append({
            'Orden': f"({p},{d},{q})({P},{D},{Q})_{S_wti}",
            'AIC': round(res.aic, 2),
            'BIC': round(res.bic, 2)
        })
        print(f"  SARIMA{(p,d,q)}{(P,D,Q)}_{S_wti} "
              f"AIC={res.aic:.2f}  BIC={res.bic:.2f}")
    except Exception as e:
        print(f"  Error en ({p},{d},{q})({P},{D},{Q}): {e}")

df_wti = pd.DataFrame(resultados_wti).sort_values('AIC')
print("\n=== Ranking por AIC ===")
print(df_wti)

# ── 7. Modelo final: SARIMA(1,1,1)(0,1,1)_52 ───────────────
modelo_wti = SARIMAX(log_precio,
                     order=(1, 1, 1),
                     seasonal_order=(0, 1, 1, 52),
                     trend='n').fit(disp=False, maxiter=400)
print("\n" + "="*60)
print(modelo_wti.summary())

# ── 8. Diagnostico ──────────────────────────────────────────
resid_wti = modelo_wti.resid
fig, axes = plt.subplots(2, 2, figsize=(12, 8))
axes[0, 0].plot(resid_wti, linewidth=0.6, color='navy')
axes[0, 0].set_title('Residuos')
plot_acf(resid_wti, lags=30, ax=axes[0, 1], title='ACF Residuos')
axes[1, 0].hist(resid_wti, bins=40, color='steelblue', edgecolor='white')
axes[1, 0].set_title('Histograma')
stats.probplot(resid_wti, dist='norm', plot=axes[1, 1])
plt.tight_layout()
plt.savefig('wti_diagnostico.png', dpi=150, bbox_inches='tight')
plt.show()

lb_wti = acorr_ljungbox(resid_wti, lags=[10, 20, 52], return_df=True)
print("\n=== Ljung-Box (WTI) ===")
print(lb_wti)

# ── 9. Prediccion 8 semanas ─────────────────────────────────
n_steps = 8
fc_wti = modelo_wti.get_forecast(steps=n_steps)
fc_log = fc_wti.predicted_mean
fc_ci_log = fc_wti.conf_int(alpha=0.05)

# Revertir transformacion logaritmica
fc_precio = np.exp(fc_log)
fc_ci_precio = np.exp(fc_ci_log)

print("\n=== Prediccion WTI (USD/barril) - 8 semanas ===")
df_fc = pd.DataFrame({
    'Pred (USD)': fc_precio.round(2),
    'IC_inf': fc_ci_precio.iloc[:, 0].round(2),
    'IC_sup': fc_ci_precio.iloc[:, 1].round(2)
})
print(df_fc)

# Grafica de prediccion
fig, ax = plt.subplots(figsize=(12, 5))
hist_plot = precio_wti[-52:]
ax.plot(hist_plot.index, hist_plot, color='saddlebrown',
        label='Precio historico (1 ano)')
ax.plot(fc_precio.index, fc_precio, 'o--',
        color='darkorange', label='Prediccion 8 semanas')
ax.fill_between(fc_ci_precio.index,
                fc_ci_precio.iloc[:, 0],
                fc_ci_precio.iloc[:, 1],
                alpha=0.3, color='orange', label='IC 95%')
ax.set_title('Prediccion precio WTI - SARIMA(1,1,1)(0,1,1)_52')
ax.set_ylabel('USD por barril')
ax.legend()
plt.tight_layout()
plt.savefig('wti_forecast.png', dpi=150, bbox_inches='tight')
plt.show()
\end{lstlisting}

\subsubsection{Interpretación de resultados}

Los precios del petróleo WTI en niveles son claramente no estacionarios ($I(1)$): la ADF no
rechaza la raíz unitaria en log-precios, pero sí en la primera diferencia. Al aplicar también
la diferencia estacional de orden 52 semanas, se elimina la autocorrelación estacional visible
en los picos de la ACF en rezagos múltiplos de 52. El modelo SARIMA$(1,1,1)(0,1,1)_{52}$ es
un punto de partida estándar (equivalente al modelo ``airline'' estacional), donde el parámetro
MA estacional $\Theta_1$ captura el ajuste ante choques anuales. La predicción se transforma
de vuelta a la escala de precios mediante $\hat{P}_t = \exp(\hat{\ell}_t)$, con intervalos
de confianza que reflejan la alta incertidumbre de los precios de commodities.

% ────────────────────────────────────────────────────────────────────────────
\section{Conclusiones}
% ────────────────────────────────────────────────────────────────────────────

\subsection{Síntesis metodológica}

Los modelos SARIMA ofrecen un marco riguroso y parsimoniosos para el análisis de series de
tiempo financieras. La metodología Box-Jenkins proporciona un procedimiento sistemático:
identificar la estructura mediante ACF y PACF, estimar por máxima verosimilitud, y validar
mediante diagnóstico de residuos. Los criterios AIC y BIC guían la selección entre modelos
competidores balanceando ajuste y complejidad.

\subsection{Aplicaciones en finanzas reales}

En la práctica financiera, los modelos SARIMA tienen aplicaciones directas en:

\begin{enumerate}
  \item \textbf{Gestión de riesgo}: la predicción de volatilidad de rendimientos (vía la
        varianza de los errores de pronóstico) alimenta modelos de VaR y CVaR.
  \item \textbf{Trading sistemático}: señales basadas en componentes AR/MA pueden identificar
        tendencias de corto plazo explotables (con las reservas del EMH).
  \item \textbf{Valoración de derivados}: la calibración de modelos de tasa de interés
        (como Vasicek o CIR) se apoya en el análisis ARIMA de la serie de tasas.
  \item \textbf{Planificación financiera corporativa}: pronósticos de ingresos, costos
        operativos e indicadores macroeconómicos con estacionalidad.
  \item \textbf{Regulación y supervisión}: bancos centrales utilizan modelos SARIMA para
        proyectar inflación, tipo de cambio y crecimiento económico.
\end{enumerate}

\subsection{Limitaciones y extensiones}

Los modelos SARIMA modelan exclusivamente la \textit{media condicional} de la serie y asumen
linealidad. Las principales limitaciones incluyen:

\begin{itemize}
  \item \textbf{Heterocedasticidad condicional}: los rendimientos financieros presentan
        \textit{clustering} de volatilidad, que requiere extensiones GARCH.
  \item \textbf{No linealidad}: regímenes de mercado (crisis vs. bonanza) violan el supuesto
        de linealidad; los modelos MS-ARIMA o de umbral (SETAR/STAR) son más apropiados.
  \item \textbf{Alta dimensión}: para múltiples series correlacionadas, los modelos VAR
        o VARMA son la extensión natural.
  \item \textbf{Machine learning}: modelos LSTM, Transformer o Prophet pueden capturar
        patrones complejos, aunque a costa de interpretabilidad.
\end{itemize}

A pesar de estas limitaciones, los modelos SARIMA permanecen como la referencia canónica
en cualquier análisis de series de tiempo univariadas, por su sólida fundamentación
teórica, eficiencia computacional e interpretabilidad de los coeficientes.

% ────────────────────────────────────────────────────────────────────────────
\section*{Referencias}
% ────────────────────────────────────────────────────────────────────────────
\addcontentsline{toc}{section}{Referencias}

\begin{enumerate}
  \item Box, G.~E.~P., Jenkins, G.~M., Reinsel, G.~C., \& Ljung, G.~M. (2015).
        \textit{Time Series Analysis: Forecasting and Control} (5th ed.). Wiley.

  \item Hamilton, J.~D. (1994). \textit{Time Series Analysis}. Princeton University Press.

  \item Tsay, R.~S. (2010). \textit{Analysis of Financial Time Series} (3rd ed.). Wiley.

  \item Brockwell, P.~J., \& Davis, R.~A. (2016). \textit{Introduction to Time Series
        and Forecasting} (3rd ed.). Springer.

  \item Hyndman, R.~J., \& Athanasopoulos, G. (2021). \textit{Forecasting: Principles
        and Practice} (3rd ed.). OTexts. \url{https://otexts.com/fpp3/}

  \item Akaike, H. (1974). A new look at the statistical model identification.
        \textit{IEEE Transactions on Automatic Control}, 19(6), 716--723.

  \item Schwarz, G. (1978). Estimating the dimension of a model.
        \textit{Annals of Statistics}, 6(2), 461--464.

  \item Ljung, G.~M., \& Box, G.~E.~P. (1978). On a measure of lack of fit in time
        series models. \textit{Biometrika}, 65(2), 297--303.

  \item Dickey, D.~A., \& Fuller, W.~A. (1979). Distribution of the estimators for
        autoregressive time series with a unit root. \textit{Journal of the American
        Statistical Association}, 74(366), 427--431.

  \item Seabold, S., \& Perktold, J. (2010). Statsmodels: Econometric and statistical
        modeling with Python. \textit{Proceedings of the 9th Python in Science Conference}.
\end{enumerate}

\end{document}
